\section{Appendix P -- Reference, Prediction, and Falsifier Digest}
% R005_APPENDIX_READER_FRAME
Reader frame. This appendix supports the main manuscript by explaining corpus role, evidence class, audit route, and limitation boundary. It is not a detached raw dump.


\label{sec:oc-core-1-3-reference-benchmark-atlas}

The experiment, falsifiability, and prediction corpus is preserved as evidence, not printed as a large raw atlas in the monograph. The public PDF gives the reading method and the boundaries; the evidence package carries the exact rows, technical files, and hashes.


\subsection{Prediction and falsifier reading rule}

A prediction row is public only when it names a formula or reconstruction rule, a source snapshot, a comparator, an uncertainty or residual, a negative control, a falsifier, and a replay hash. Older prediction language that lacks those fields remains an archival route and cannot promote a release claim.


\subsection{Benchmark evidence route}

The benchmark rows are useful because they show how a claim would be checked. They do not imply that every domain has been completed. The corpus register records the full benchmark tree so reviewers can audit coverage without turning the monograph into an index dump.


\subsection{Reference, falsifier, and prediction scholarly synthesis}

The reference scholarly synthesis replaces the earlier benchmark atlas dump. It keeps the same experiment, falsifier, and prediction corpus visible while forcing each route to state what it can support, how it would be checked, and what would reopen it.


\subsubsection{Experiments}

\paragraph{Experiments: Experiments Master}

For Experiments: Experiments Master, the scholarly synthesis contributes experiment design and operational test route, explains why this module exists, shows how it connects to the current 1.3.3 argument, and prevents stronger inherited language from being promoted without the current proof, replay, comparator, and falsifier surfaces.


For Experiments: Experiments Master, reader use is to ask whether the test has input, observable, comparator, expected output, and failure interpretation; if that module-specific question cannot be answered from the promoted theorem, proof sheet, finite semantic case, replay table, comparator row, or claim-boundary section, the source remains background support rather than public theorem evidence.


For Experiments: Experiments Master, the audit route keeps the complete technical module indexed by the release corpus index and the public evidence package. The science is therefore not discarded, while the PDF avoids becoming a unintegrated technical dossier. A line-level reviewer follows the register; a scientific reader follows this digest and the main chapter.


Experiments: Experiments Master reader checkpoint is satisfied only when the main manuscript tells the reader which claim class this module can support, which evidence class limits it, which prior-art or comparator boundary applies, and which failure would reopen the claim. This fourth check is deliberately positive: the module must contribute intelligibility, not merely avoid forbidden wording.


For Experiments: Experiments Master, the external wording boundary allows public language to say that the route helps organize, test, or constrain a bounded OC claim. It may not say that the route proves full-domain closure, final bounded cross-domain synthesis status, unrestricted comparative claim, or a domain result absent from the current proof and evidence layer. This boundary is part of the source's scientific content, not a marketing caveat.


\paragraph{Experiments: Experiments K0}

For Experiments: Experiments K0, the scholarly synthesis contributes experiment design and operational test route, explains why this module exists, shows how it connects to the current 1.3.3 argument, and prevents stronger inherited language from being promoted without the current proof, replay, comparator, and falsifier surfaces.


For Experiments: Experiments K0, reader use is to ask whether the test has input, observable, comparator, expected output, and failure interpretation; if that module-specific question cannot be answered from the promoted theorem, proof sheet, finite semantic case, replay table, comparator row, or claim-boundary section, the source remains background support rather than public theorem evidence.


For Experiments: Experiments K0, the audit route keeps the complete technical module indexed by the release corpus index and the public evidence package. The science is therefore not discarded, while the PDF avoids becoming a unintegrated technical dossier. A line-level reviewer follows the register; a scientific reader follows this digest and the main chapter.


Experiments: Experiments K0 reader checkpoint is satisfied only when the main manuscript tells the reader which claim class this module can support, which evidence class limits it, which prior-art or comparator boundary applies, and which failure would reopen the claim. This fourth check is deliberately positive: the module must contribute intelligibility, not merely avoid forbidden wording.


For Experiments: Experiments K0, the external wording boundary allows public language to say that the route helps organize, test, or constrain a bounded OC claim. It may not say that the route proves full-domain closure, final bounded cross-domain synthesis status, unrestricted comparative claim, or a domain result absent from the current proof and evidence layer. This boundary is part of the source's scientific content, not a marketing caveat.


\paragraph{Experiments: Experiments K1}

For Experiments: Experiments K1, the scholarly synthesis contributes experiment design and operational test route, explains why this module exists, shows how it connects to the current 1.3.3 argument, and prevents stronger inherited language from being promoted without the current proof, replay, comparator, and falsifier surfaces.


For Experiments: Experiments K1, reader use is to ask whether the test has input, observable, comparator, expected output, and failure interpretation; if that module-specific question cannot be answered from the promoted theorem, proof sheet, finite semantic case, replay table, comparator row, or claim-boundary section, the source remains background support rather than public theorem evidence.


For Experiments: Experiments K1, the audit route keeps the complete technical module indexed by the release corpus index and the public evidence package. The science is therefore not discarded, while the PDF avoids becoming a unintegrated technical dossier. A line-level reviewer follows the register; a scientific reader follows this digest and the main chapter.


Experiments: Experiments K1 reader checkpoint is satisfied only when the main manuscript tells the reader which claim class this module can support, which evidence class limits it, which prior-art or comparator boundary applies, and which failure would reopen the claim. This fourth check is deliberately positive: the module must contribute intelligibility, not merely avoid forbidden wording.


For Experiments: Experiments K1, the external wording boundary allows public language to say that the route helps organize, test, or constrain a bounded OC claim. It may not say that the route proves full-domain closure, final bounded cross-domain synthesis status, unrestricted comparative claim, or a domain result absent from the current proof and evidence layer. This boundary is part of the source's scientific content, not a marketing caveat.


\paragraph{Experiments: Experiments K2}

For Experiments: Experiments K2, the scholarly synthesis contributes experiment design and operational test route, explains why this module exists, shows how it connects to the current 1.3.3 argument, and prevents stronger inherited language from being promoted without the current proof, replay, comparator, and falsifier surfaces.


For Experiments: Experiments K2, reader use is to ask whether the test has input, observable, comparator, expected output, and failure interpretation; if that module-specific question cannot be answered from the promoted theorem, proof sheet, finite semantic case, replay table, comparator row, or claim-boundary section, the source remains background support rather than public theorem evidence.


For Experiments: Experiments K2, the audit route keeps the complete technical module indexed by the release corpus index and the public evidence package. The science is therefore not discarded, while the PDF avoids becoming a unintegrated technical dossier. A line-level reviewer follows the register; a scientific reader follows this digest and the main chapter.


Experiments: Experiments K2 reader checkpoint is satisfied only when the main manuscript tells the reader which claim class this module can support, which evidence class limits it, which prior-art or comparator boundary applies, and which failure would reopen the claim. This fourth check is deliberately positive: the module must contribute intelligibility, not merely avoid forbidden wording.


For Experiments: Experiments K2, the external wording boundary allows public language to say that the route helps organize, test, or constrain a bounded OC claim. It may not say that the route proves full-domain closure, final bounded cross-domain synthesis status, unrestricted comparative claim, or a domain result absent from the current proof and evidence layer. This boundary is part of the source's scientific content, not a marketing caveat.


\paragraph{Experiments: Experiments K3}

For Experiments: Experiments K3, the scholarly synthesis contributes experiment design and operational test route, explains why this module exists, shows how it connects to the current 1.3.3 argument, and prevents stronger inherited language from being promoted without the current proof, replay, comparator, and falsifier surfaces.


For Experiments: Experiments K3, reader use is to ask whether the test has input, observable, comparator, expected output, and failure interpretation; if that module-specific question cannot be answered from the promoted theorem, proof sheet, finite semantic case, replay table, comparator row, or claim-boundary section, the source remains background support rather than public theorem evidence.


For Experiments: Experiments K3, the audit route keeps the complete technical module indexed by the release corpus index and the public evidence package. The science is therefore not discarded, while the PDF avoids becoming a unintegrated technical dossier. A line-level reviewer follows the register; a scientific reader follows this digest and the main chapter.


Experiments: Experiments K3 reader checkpoint is satisfied only when the main manuscript tells the reader which claim class this module can support, which evidence class limits it, which prior-art or comparator boundary applies, and which failure would reopen the claim. This fourth check is deliberately positive: the module must contribute intelligibility, not merely avoid forbidden wording.


For Experiments: Experiments K3, the external wording boundary allows public language to say that the route helps organize, test, or constrain a bounded OC claim. It may not say that the route proves full-domain closure, final bounded cross-domain synthesis status, unrestricted comparative claim, or a domain result absent from the current proof and evidence layer. This boundary is part of the source's scientific content, not a marketing caveat.


\paragraph{Experiments: Experiments K4}

For Experiments: Experiments K4, the scholarly synthesis contributes experiment design and operational test route, explains why this module exists, shows how it connects to the current 1.3.3 argument, and prevents stronger inherited language from being promoted without the current proof, replay, comparator, and falsifier surfaces.


For Experiments: Experiments K4, reader use is to ask whether the test has input, observable, comparator, expected output, and failure interpretation; if that module-specific question cannot be answered from the promoted theorem, proof sheet, finite semantic case, replay table, comparator row, or claim-boundary section, the source remains background support rather than public theorem evidence.


For Experiments: Experiments K4, the audit route keeps the complete technical module indexed by the release corpus index and the public evidence package. The science is therefore not discarded, while the PDF avoids becoming a unintegrated technical dossier. A line-level reviewer follows the register; a scientific reader follows this digest and the main chapter.


Experiments: Experiments K4 reader checkpoint is satisfied only when the main manuscript tells the reader which claim class this module can support, which evidence class limits it, which prior-art or comparator boundary applies, and which failure would reopen the claim. This fourth check is deliberately positive: the module must contribute intelligibility, not merely avoid forbidden wording.


For Experiments: Experiments K4, the external wording boundary allows public language to say that the route helps organize, test, or constrain a bounded OC claim. It may not say that the route proves full-domain closure, final bounded cross-domain synthesis status, unrestricted comparative claim, or a domain result absent from the current proof and evidence layer. This boundary is part of the source's scientific content, not a marketing caveat.


\paragraph{Experiments: Experiments K5}

For Experiments: Experiments K5, the scholarly synthesis contributes experiment design and operational test route, explains why this module exists, shows how it connects to the current 1.3.3 argument, and prevents stronger inherited language from being promoted without the current proof, replay, comparator, and falsifier surfaces.


For Experiments: Experiments K5, reader use is to ask whether the test has input, observable, comparator, expected output, and failure interpretation; if that module-specific question cannot be answered from the promoted theorem, proof sheet, finite semantic case, replay table, comparator row, or claim-boundary section, the source remains background support rather than public theorem evidence.


For Experiments: Experiments K5, the audit route keeps the complete technical module indexed by the release corpus index and the public evidence package. The science is therefore not discarded, while the PDF avoids becoming a unintegrated technical dossier. A line-level reviewer follows the register; a scientific reader follows this digest and the main chapter.


Experiments: Experiments K5 reader checkpoint is satisfied only when the main manuscript tells the reader which claim class this module can support, which evidence class limits it, which prior-art or comparator boundary applies, and which failure would reopen the claim. This fourth check is deliberately positive: the module must contribute intelligibility, not merely avoid forbidden wording.


For Experiments: Experiments K5, the external wording boundary allows public language to say that the route helps organize, test, or constrain a bounded OC claim. It may not say that the route proves full-domain closure, final bounded cross-domain synthesis status, unrestricted comparative claim, or a domain result absent from the current proof and evidence layer. This boundary is part of the source's scientific content, not a marketing caveat.


\paragraph{Experiments: Experiments K6}

For Experiments: Experiments K6, the scholarly synthesis contributes experiment design and operational test route, explains why this module exists, shows how it connects to the current 1.3.3 argument, and prevents stronger inherited language from being promoted without the current proof, replay, comparator, and falsifier surfaces.


For Experiments: Experiments K6, reader use is to ask whether the test has input, observable, comparator, expected output, and failure interpretation; if that module-specific question cannot be answered from the promoted theorem, proof sheet, finite semantic case, replay table, comparator row, or claim-boundary section, the source remains background support rather than public theorem evidence.


For Experiments: Experiments K6, the audit route keeps the complete technical module indexed by the release corpus index and the public evidence package. The science is therefore not discarded, while the PDF avoids becoming a unintegrated technical dossier. A line-level reviewer follows the register; a scientific reader follows this digest and the main chapter.


Experiments: Experiments K6 reader checkpoint is satisfied only when the main manuscript tells the reader which claim class this module can support, which evidence class limits it, which prior-art or comparator boundary applies, and which failure would reopen the claim. This fourth check is deliberately positive: the module must contribute intelligibility, not merely avoid forbidden wording.


For Experiments: Experiments K6, the external wording boundary allows public language to say that the route helps organize, test, or constrain a bounded OC claim. It may not say that the route proves full-domain closure, final bounded cross-domain synthesis status, unrestricted comparative claim, or a domain result absent from the current proof and evidence layer. This boundary is part of the source's scientific content, not a marketing caveat.


\paragraph{Experiments: Experiments K7}

For Experiments: Experiments K7, the scholarly synthesis contributes experiment design and operational test route, explains why this module exists, shows how it connects to the current 1.3.3 argument, and prevents stronger inherited language from being promoted without the current proof, replay, comparator, and falsifier surfaces.


For Experiments: Experiments K7, reader use is to ask whether the test has input, observable, comparator, expected output, and failure interpretation; if that module-specific question cannot be answered from the promoted theorem, proof sheet, finite semantic case, replay table, comparator row, or claim-boundary section, the source remains background support rather than public theorem evidence.


For Experiments: Experiments K7, the audit route keeps the complete technical module indexed by the release corpus index and the public evidence package. The science is therefore not discarded, while the PDF avoids becoming a unintegrated technical dossier. A line-level reviewer follows the register; a scientific reader follows this digest and the main chapter.


Experiments: Experiments K7 reader checkpoint is satisfied only when the main manuscript tells the reader which claim class this module can support, which evidence class limits it, which prior-art or comparator boundary applies, and which failure would reopen the claim. This fourth check is deliberately positive: the module must contribute intelligibility, not merely avoid forbidden wording.


For Experiments: Experiments K7, the external wording boundary allows public language to say that the route helps organize, test, or constrain a bounded OC claim. It may not say that the route proves full-domain closure, final bounded cross-domain synthesis status, unrestricted comparative claim, or a domain result absent from the current proof and evidence layer. This boundary is part of the source's scientific content, not a marketing caveat.


\paragraph{Experiments: Experiments K8}

For Experiments: Experiments K8, the scholarly synthesis contributes experiment design and operational test route, explains why this module exists, shows how it connects to the current 1.3.3 argument, and prevents stronger inherited language from being promoted without the current proof, replay, comparator, and falsifier surfaces.


For Experiments: Experiments K8, reader use is to ask whether the test has input, observable, comparator, expected output, and failure interpretation; if that module-specific question cannot be answered from the promoted theorem, proof sheet, finite semantic case, replay table, comparator row, or claim-boundary section, the source remains background support rather than public theorem evidence.


For Experiments: Experiments K8, the audit route keeps the complete technical module indexed by the release corpus index and the public evidence package. The science is therefore not discarded, while the PDF avoids becoming a unintegrated technical dossier. A line-level reviewer follows the register; a scientific reader follows this digest and the main chapter.


Experiments: Experiments K8 reader checkpoint is satisfied only when the main manuscript tells the reader which claim class this module can support, which evidence class limits it, which prior-art or comparator boundary applies, and which failure would reopen the claim. This fourth check is deliberately positive: the module must contribute intelligibility, not merely avoid forbidden wording.


For Experiments: Experiments K8, the external wording boundary allows public language to say that the route helps organize, test, or constrain a bounded OC claim. It may not say that the route proves full-domain closure, final bounded cross-domain synthesis status, unrestricted comparative claim, or a domain result absent from the current proof and evidence layer. This boundary is part of the source's scientific content, not a marketing caveat.


\paragraph{Experiments: Experiments K9}

For Experiments: Experiments K9, the scholarly synthesis contributes experiment design and operational test route, explains why this module exists, shows how it connects to the current 1.3.3 argument, and prevents stronger inherited language from being promoted without the current proof, replay, comparator, and falsifier surfaces.


For Experiments: Experiments K9, reader use is to ask whether the test has input, observable, comparator, expected output, and failure interpretation; if that module-specific question cannot be answered from the promoted theorem, proof sheet, finite semantic case, replay table, comparator row, or claim-boundary section, the source remains background support rather than public theorem evidence.


For Experiments: Experiments K9, the audit route keeps the complete technical module indexed by the release corpus index and the public evidence package. The science is therefore not discarded, while the PDF avoids becoming a unintegrated technical dossier. A line-level reviewer follows the register; a scientific reader follows this digest and the main chapter.


Experiments: Experiments K9 reader checkpoint is satisfied only when the main manuscript tells the reader which claim class this module can support, which evidence class limits it, which prior-art or comparator boundary applies, and which failure would reopen the claim. This fourth check is deliberately positive: the module must contribute intelligibility, not merely avoid forbidden wording.


For Experiments: Experiments K9, the external wording boundary allows public language to say that the route helps organize, test, or constrain a bounded OC claim. It may not say that the route proves full-domain closure, final bounded cross-domain synthesis status, unrestricted comparative claim, or a domain result absent from the current proof and evidence layer. This boundary is part of the source's scientific content, not a marketing caveat.


\paragraph{Experiments: Experiments K10}

For Experiments: Experiments K10, the scholarly synthesis contributes experiment design and operational test route, explains why this module exists, shows how it connects to the current 1.3.3 argument, and prevents stronger inherited language from being promoted without the current proof, replay, comparator, and falsifier surfaces.


For Experiments: Experiments K10, reader use is to ask whether the test has input, observable, comparator, expected output, and failure interpretation; if that module-specific question cannot be answered from the promoted theorem, proof sheet, finite semantic case, replay table, comparator row, or claim-boundary section, the source remains background support rather than public theorem evidence.


For Experiments: Experiments K10, the audit route keeps the complete technical module indexed by the release corpus index and the public evidence package. The science is therefore not discarded, while the PDF avoids becoming a unintegrated technical dossier. A line-level reviewer follows the register; a scientific reader follows this digest and the main chapter.


Experiments: Experiments K10 reader checkpoint is satisfied only when the main manuscript tells the reader which claim class this module can support, which evidence class limits it, which prior-art or comparator boundary applies, and which failure would reopen the claim. This fourth check is deliberately positive: the module must contribute intelligibility, not merely avoid forbidden wording.


For Experiments: Experiments K10, the external wording boundary allows public language to say that the route helps organize, test, or constrain a bounded OC claim. It may not say that the route proves full-domain closure, final bounded cross-domain synthesis status, unrestricted comparative claim, or a domain result absent from the current proof and evidence layer. This boundary is part of the source's scientific content, not a marketing caveat.


Integration checkpoint 12. The preceding modules teach inspection logic; they do not inflate claim strength. The release remains bounded by the current proof and evidence surface.


\paragraph{Experiments: Experiments K11}

For Experiments: Experiments K11, the scholarly synthesis contributes experiment design and operational test route, explains why this module exists, shows how it connects to the current 1.3.3 argument, and prevents stronger inherited language from being promoted without the current proof, replay, comparator, and falsifier surfaces.


For Experiments: Experiments K11, reader use is to ask whether the test has input, observable, comparator, expected output, and failure interpretation; if that module-specific question cannot be answered from the promoted theorem, proof sheet, finite semantic case, replay table, comparator row, or claim-boundary section, the source remains background support rather than public theorem evidence.


For Experiments: Experiments K11, the audit route keeps the complete technical module indexed by the release corpus index and the public evidence package. The science is therefore not discarded, while the PDF avoids becoming a unintegrated technical dossier. A line-level reviewer follows the register; a scientific reader follows this digest and the main chapter.


Experiments: Experiments K11 reader checkpoint is satisfied only when the main manuscript tells the reader which claim class this module can support, which evidence class limits it, which prior-art or comparator boundary applies, and which failure would reopen the claim. This fourth check is deliberately positive: the module must contribute intelligibility, not merely avoid forbidden wording.


For Experiments: Experiments K11, the external wording boundary allows public language to say that the route helps organize, test, or constrain a bounded OC claim. It may not say that the route proves full-domain closure, final bounded cross-domain synthesis status, unrestricted comparative claim, or a domain result absent from the current proof and evidence layer. This boundary is part of the source's scientific content, not a marketing caveat.


\paragraph{Experiments: Experiments K12}

For Experiments: Experiments K12, the scholarly synthesis contributes experiment design and operational test route, explains why this module exists, shows how it connects to the current 1.3.3 argument, and prevents stronger inherited language from being promoted without the current proof, replay, comparator, and falsifier surfaces.


For Experiments: Experiments K12, reader use is to ask whether the test has input, observable, comparator, expected output, and failure interpretation; if that module-specific question cannot be answered from the promoted theorem, proof sheet, finite semantic case, replay table, comparator row, or claim-boundary section, the source remains background support rather than public theorem evidence.


For Experiments: Experiments K12, the audit route keeps the complete technical module indexed by the release corpus index and the public evidence package. The science is therefore not discarded, while the PDF avoids becoming a unintegrated technical dossier. A line-level reviewer follows the register; a scientific reader follows this digest and the main chapter.


Experiments: Experiments K12 reader checkpoint is satisfied only when the main manuscript tells the reader which claim class this module can support, which evidence class limits it, which prior-art or comparator boundary applies, and which failure would reopen the claim. This fourth check is deliberately positive: the module must contribute intelligibility, not merely avoid forbidden wording.


For Experiments: Experiments K12, the external wording boundary allows public language to say that the route helps organize, test, or constrain a bounded OC claim. It may not say that the route proves full-domain closure, final bounded cross-domain synthesis status, unrestricted comparative claim, or a domain result absent from the current proof and evidence layer. This boundary is part of the source's scientific content, not a marketing caveat.


\subsubsection{Falsifiability}

\paragraph{Falsifiability: Falsifiability Master}

For Falsifiability: Falsifiability Master, the scholarly synthesis contributes falsifier grammar and reopening condition, explains why this module exists, shows how it connects to the current 1.3.3 argument, and prevents stronger inherited language from being promoted without the current proof, replay, comparator, and falsifier surfaces.


For Falsifiability: Falsifiability Master, reader use is to ask whether the claimed phenomenon would actually reopen under the stated falsifier; if that module-specific question cannot be answered from the promoted theorem, proof sheet, finite semantic case, replay table, comparator row, or claim-boundary section, the source remains background support rather than public theorem evidence.


For Falsifiability: Falsifiability Master, the audit route keeps the complete technical module indexed by the release corpus index and the public evidence package. The science is therefore not discarded, while the PDF avoids becoming a unintegrated technical dossier. A line-level reviewer follows the register; a scientific reader follows this digest and the main chapter.


Falsifiability: Falsifiability Master reader checkpoint is satisfied only when the main manuscript tells the reader which claim class this module can support, which evidence class limits it, which prior-art or comparator boundary applies, and which failure would reopen the claim. This fourth check is deliberately positive: the module must contribute intelligibility, not merely avoid forbidden wording.


For Falsifiability: Falsifiability Master, the external wording boundary allows public language to say that the route helps organize, test, or constrain a bounded OC claim. It may not say that the route proves full-domain closure, final bounded cross-domain synthesis status, unrestricted comparative claim, or a domain result absent from the current proof and evidence layer. This boundary is part of the source's scientific content, not a marketing caveat.


\paragraph{Falsifiability: Falsifiability K0}

For Falsifiability: Falsifiability K0, the scholarly synthesis contributes falsifier grammar and reopening condition, explains why this module exists, shows how it connects to the current 1.3.3 argument, and prevents stronger inherited language from being promoted without the current proof, replay, comparator, and falsifier surfaces.


For Falsifiability: Falsifiability K0, reader use is to ask whether the claimed phenomenon would actually reopen under the stated falsifier; if that module-specific question cannot be answered from the promoted theorem, proof sheet, finite semantic case, replay table, comparator row, or claim-boundary section, the source remains background support rather than public theorem evidence.


For Falsifiability: Falsifiability K0, the audit route keeps the complete technical module indexed by the release corpus index and the public evidence package. The science is therefore not discarded, while the PDF avoids becoming a unintegrated technical dossier. A line-level reviewer follows the register; a scientific reader follows this digest and the main chapter.


Falsifiability: Falsifiability K0 reader checkpoint is satisfied only when the main manuscript tells the reader which claim class this module can support, which evidence class limits it, which prior-art or comparator boundary applies, and which failure would reopen the claim. This fourth check is deliberately positive: the module must contribute intelligibility, not merely avoid forbidden wording.


For Falsifiability: Falsifiability K0, the external wording boundary allows public language to say that the route helps organize, test, or constrain a bounded OC claim. It may not say that the route proves full-domain closure, final bounded cross-domain synthesis status, unrestricted comparative claim, or a domain result absent from the current proof and evidence layer. This boundary is part of the source's scientific content, not a marketing caveat.


\paragraph{Falsifiability: Falsifiability K1}

For Falsifiability: Falsifiability K1, the scholarly synthesis contributes falsifier grammar and reopening condition, explains why this module exists, shows how it connects to the current 1.3.3 argument, and prevents stronger inherited language from being promoted without the current proof, replay, comparator, and falsifier surfaces.


For Falsifiability: Falsifiability K1, reader use is to ask whether the claimed phenomenon would actually reopen under the stated falsifier; if that module-specific question cannot be answered from the promoted theorem, proof sheet, finite semantic case, replay table, comparator row, or claim-boundary section, the source remains background support rather than public theorem evidence.


For Falsifiability: Falsifiability K1, the audit route keeps the complete technical module indexed by the release corpus index and the public evidence package. The science is therefore not discarded, while the PDF avoids becoming a unintegrated technical dossier. A line-level reviewer follows the register; a scientific reader follows this digest and the main chapter.


Falsifiability: Falsifiability K1 reader checkpoint is satisfied only when the main manuscript tells the reader which claim class this module can support, which evidence class limits it, which prior-art or comparator boundary applies, and which failure would reopen the claim. This fourth check is deliberately positive: the module must contribute intelligibility, not merely avoid forbidden wording.


For Falsifiability: Falsifiability K1, the external wording boundary allows public language to say that the route helps organize, test, or constrain a bounded OC claim. It may not say that the route proves full-domain closure, final bounded cross-domain synthesis status, unrestricted comparative claim, or a domain result absent from the current proof and evidence layer. This boundary is part of the source's scientific content, not a marketing caveat.


\paragraph{Falsifiability: Falsifiability K2}

For Falsifiability: Falsifiability K2, the scholarly synthesis contributes falsifier grammar and reopening condition, explains why this module exists, shows how it connects to the current 1.3.3 argument, and prevents stronger inherited language from being promoted without the current proof, replay, comparator, and falsifier surfaces.


For Falsifiability: Falsifiability K2, reader use is to ask whether the claimed phenomenon would actually reopen under the stated falsifier; if that module-specific question cannot be answered from the promoted theorem, proof sheet, finite semantic case, replay table, comparator row, or claim-boundary section, the source remains background support rather than public theorem evidence.


For Falsifiability: Falsifiability K2, the audit route keeps the complete technical module indexed by the release corpus index and the public evidence package. The science is therefore not discarded, while the PDF avoids becoming a unintegrated technical dossier. A line-level reviewer follows the register; a scientific reader follows this digest and the main chapter.


Falsifiability: Falsifiability K2 reader checkpoint is satisfied only when the main manuscript tells the reader which claim class this module can support, which evidence class limits it, which prior-art or comparator boundary applies, and which failure would reopen the claim. This fourth check is deliberately positive: the module must contribute intelligibility, not merely avoid forbidden wording.


For Falsifiability: Falsifiability K2, the external wording boundary allows public language to say that the route helps organize, test, or constrain a bounded OC claim. It may not say that the route proves full-domain closure, final bounded cross-domain synthesis status, unrestricted comparative claim, or a domain result absent from the current proof and evidence layer. This boundary is part of the source's scientific content, not a marketing caveat.


\paragraph{Falsifiability: Falsifiability K3}

For Falsifiability: Falsifiability K3, the scholarly synthesis contributes falsifier grammar and reopening condition, explains why this module exists, shows how it connects to the current 1.3.3 argument, and prevents stronger inherited language from being promoted without the current proof, replay, comparator, and falsifier surfaces.


For Falsifiability: Falsifiability K3, reader use is to ask whether the claimed phenomenon would actually reopen under the stated falsifier; if that module-specific question cannot be answered from the promoted theorem, proof sheet, finite semantic case, replay table, comparator row, or claim-boundary section, the source remains background support rather than public theorem evidence.


For Falsifiability: Falsifiability K3, the audit route keeps the complete technical module indexed by the release corpus index and the public evidence package. The science is therefore not discarded, while the PDF avoids becoming a unintegrated technical dossier. A line-level reviewer follows the register; a scientific reader follows this digest and the main chapter.


Falsifiability: Falsifiability K3 reader checkpoint is satisfied only when the main manuscript tells the reader which claim class this module can support, which evidence class limits it, which prior-art or comparator boundary applies, and which failure would reopen the claim. This fourth check is deliberately positive: the module must contribute intelligibility, not merely avoid forbidden wording.


For Falsifiability: Falsifiability K3, the external wording boundary allows public language to say that the route helps organize, test, or constrain a bounded OC claim. It may not say that the route proves full-domain closure, final bounded cross-domain synthesis status, unrestricted comparative claim, or a domain result absent from the current proof and evidence layer. This boundary is part of the source's scientific content, not a marketing caveat.


\paragraph{Falsifiability: Falsifiability K4}

For Falsifiability: Falsifiability K4, the scholarly synthesis contributes falsifier grammar and reopening condition, explains why this module exists, shows how it connects to the current 1.3.3 argument, and prevents stronger inherited language from being promoted without the current proof, replay, comparator, and falsifier surfaces.


For Falsifiability: Falsifiability K4, reader use is to ask whether the claimed phenomenon would actually reopen under the stated falsifier; if that module-specific question cannot be answered from the promoted theorem, proof sheet, finite semantic case, replay table, comparator row, or claim-boundary section, the source remains background support rather than public theorem evidence.


For Falsifiability: Falsifiability K4, the audit route keeps the complete technical module indexed by the release corpus index and the public evidence package. The science is therefore not discarded, while the PDF avoids becoming a unintegrated technical dossier. A line-level reviewer follows the register; a scientific reader follows this digest and the main chapter.


Falsifiability: Falsifiability K4 reader checkpoint is satisfied only when the main manuscript tells the reader which claim class this module can support, which evidence class limits it, which prior-art or comparator boundary applies, and which failure would reopen the claim. This fourth check is deliberately positive: the module must contribute intelligibility, not merely avoid forbidden wording.


For Falsifiability: Falsifiability K4, the external wording boundary allows public language to say that the route helps organize, test, or constrain a bounded OC claim. It may not say that the route proves full-domain closure, final bounded cross-domain synthesis status, unrestricted comparative claim, or a domain result absent from the current proof and evidence layer. This boundary is part of the source's scientific content, not a marketing caveat.


\paragraph{Falsifiability: Falsifiability K5}

For Falsifiability: Falsifiability K5, the scholarly synthesis contributes falsifier grammar and reopening condition, explains why this module exists, shows how it connects to the current 1.3.3 argument, and prevents stronger inherited language from being promoted without the current proof, replay, comparator, and falsifier surfaces.


For Falsifiability: Falsifiability K5, reader use is to ask whether the claimed phenomenon would actually reopen under the stated falsifier; if that module-specific question cannot be answered from the promoted theorem, proof sheet, finite semantic case, replay table, comparator row, or claim-boundary section, the source remains background support rather than public theorem evidence.


For Falsifiability: Falsifiability K5, the audit route keeps the complete technical module indexed by the release corpus index and the public evidence package. The science is therefore not discarded, while the PDF avoids becoming a unintegrated technical dossier. A line-level reviewer follows the register; a scientific reader follows this digest and the main chapter.


Falsifiability: Falsifiability K5 reader checkpoint is satisfied only when the main manuscript tells the reader which claim class this module can support, which evidence class limits it, which prior-art or comparator boundary applies, and which failure would reopen the claim. This fourth check is deliberately positive: the module must contribute intelligibility, not merely avoid forbidden wording.


For Falsifiability: Falsifiability K5, the external wording boundary allows public language to say that the route helps organize, test, or constrain a bounded OC claim. It may not say that the route proves full-domain closure, final bounded cross-domain synthesis status, unrestricted comparative claim, or a domain result absent from the current proof and evidence layer. This boundary is part of the source's scientific content, not a marketing caveat.


\paragraph{Falsifiability: Falsifiability K6}

For Falsifiability: Falsifiability K6, the scholarly synthesis contributes falsifier grammar and reopening condition, explains why this module exists, shows how it connects to the current 1.3.3 argument, and prevents stronger inherited language from being promoted without the current proof, replay, comparator, and falsifier surfaces.


For Falsifiability: Falsifiability K6, reader use is to ask whether the claimed phenomenon would actually reopen under the stated falsifier; if that module-specific question cannot be answered from the promoted theorem, proof sheet, finite semantic case, replay table, comparator row, or claim-boundary section, the source remains background support rather than public theorem evidence.


For Falsifiability: Falsifiability K6, the audit route keeps the complete technical module indexed by the release corpus index and the public evidence package. The science is therefore not discarded, while the PDF avoids becoming a unintegrated technical dossier. A line-level reviewer follows the register; a scientific reader follows this digest and the main chapter.


Falsifiability: Falsifiability K6 reader checkpoint is satisfied only when the main manuscript tells the reader which claim class this module can support, which evidence class limits it, which prior-art or comparator boundary applies, and which failure would reopen the claim. This fourth check is deliberately positive: the module must contribute intelligibility, not merely avoid forbidden wording.


For Falsifiability: Falsifiability K6, the external wording boundary allows public language to say that the route helps organize, test, or constrain a bounded OC claim. It may not say that the route proves full-domain closure, final bounded cross-domain synthesis status, unrestricted comparative claim, or a domain result absent from the current proof and evidence layer. This boundary is part of the source's scientific content, not a marketing caveat.


\paragraph{Falsifiability: Falsifiability K7}

For Falsifiability: Falsifiability K7, the scholarly synthesis contributes falsifier grammar and reopening condition, explains why this module exists, shows how it connects to the current 1.3.3 argument, and prevents stronger inherited language from being promoted without the current proof, replay, comparator, and falsifier surfaces.


For Falsifiability: Falsifiability K7, reader use is to ask whether the claimed phenomenon would actually reopen under the stated falsifier; if that module-specific question cannot be answered from the promoted theorem, proof sheet, finite semantic case, replay table, comparator row, or claim-boundary section, the source remains background support rather than public theorem evidence.


For Falsifiability: Falsifiability K7, the audit route keeps the complete technical module indexed by the release corpus index and the public evidence package. The science is therefore not discarded, while the PDF avoids becoming a unintegrated technical dossier. A line-level reviewer follows the register; a scientific reader follows this digest and the main chapter.


Falsifiability: Falsifiability K7 reader checkpoint is satisfied only when the main manuscript tells the reader which claim class this module can support, which evidence class limits it, which prior-art or comparator boundary applies, and which failure would reopen the claim. This fourth check is deliberately positive: the module must contribute intelligibility, not merely avoid forbidden wording.


For Falsifiability: Falsifiability K7, the external wording boundary allows public language to say that the route helps organize, test, or constrain a bounded OC claim. It may not say that the route proves full-domain closure, final bounded cross-domain synthesis status, unrestricted comparative claim, or a domain result absent from the current proof and evidence layer. This boundary is part of the source's scientific content, not a marketing caveat.


\paragraph{Falsifiability: Falsifiability K8}

For Falsifiability: Falsifiability K8, the scholarly synthesis contributes falsifier grammar and reopening condition, explains why this module exists, shows how it connects to the current 1.3.3 argument, and prevents stronger inherited language from being promoted without the current proof, replay, comparator, and falsifier surfaces.


For Falsifiability: Falsifiability K8, reader use is to ask whether the claimed phenomenon would actually reopen under the stated falsifier; if that module-specific question cannot be answered from the promoted theorem, proof sheet, finite semantic case, replay table, comparator row, or claim-boundary section, the source remains background support rather than public theorem evidence.


For Falsifiability: Falsifiability K8, the audit route keeps the complete technical module indexed by the release corpus index and the public evidence package. The science is therefore not discarded, while the PDF avoids becoming a unintegrated technical dossier. A line-level reviewer follows the register; a scientific reader follows this digest and the main chapter.


Falsifiability: Falsifiability K8 reader checkpoint is satisfied only when the main manuscript tells the reader which claim class this module can support, which evidence class limits it, which prior-art or comparator boundary applies, and which failure would reopen the claim. This fourth check is deliberately positive: the module must contribute intelligibility, not merely avoid forbidden wording.


For Falsifiability: Falsifiability K8, the external wording boundary allows public language to say that the route helps organize, test, or constrain a bounded OC claim. It may not say that the route proves full-domain closure, final bounded cross-domain synthesis status, unrestricted comparative claim, or a domain result absent from the current proof and evidence layer. This boundary is part of the source's scientific content, not a marketing caveat.


Integration checkpoint 24. The preceding modules teach inspection logic; they do not inflate claim strength. The release remains bounded by the current proof and evidence surface.


\paragraph{Falsifiability: Falsifiability K9}

For Falsifiability: Falsifiability K9, the scholarly synthesis contributes falsifier grammar and reopening condition, explains why this module exists, shows how it connects to the current 1.3.3 argument, and prevents stronger inherited language from being promoted without the current proof, replay, comparator, and falsifier surfaces.


For Falsifiability: Falsifiability K9, reader use is to ask whether the claimed phenomenon would actually reopen under the stated falsifier; if that module-specific question cannot be answered from the promoted theorem, proof sheet, finite semantic case, replay table, comparator row, or claim-boundary section, the source remains background support rather than public theorem evidence.


For Falsifiability: Falsifiability K9, the audit route keeps the complete technical module indexed by the release corpus index and the public evidence package. The science is therefore not discarded, while the PDF avoids becoming a unintegrated technical dossier. A line-level reviewer follows the register; a scientific reader follows this digest and the main chapter.


Falsifiability: Falsifiability K9 reader checkpoint is satisfied only when the main manuscript tells the reader which claim class this module can support, which evidence class limits it, which prior-art or comparator boundary applies, and which failure would reopen the claim. This fourth check is deliberately positive: the module must contribute intelligibility, not merely avoid forbidden wording.


For Falsifiability: Falsifiability K9, the external wording boundary allows public language to say that the route helps organize, test, or constrain a bounded OC claim. It may not say that the route proves full-domain closure, final bounded cross-domain synthesis status, unrestricted comparative claim, or a domain result absent from the current proof and evidence layer. This boundary is part of the source's scientific content, not a marketing caveat.


\paragraph{Falsifiability: Falsifiability K10}

For Falsifiability: Falsifiability K10, the scholarly synthesis contributes falsifier grammar and reopening condition, explains why this module exists, shows how it connects to the current 1.3.3 argument, and prevents stronger inherited language from being promoted without the current proof, replay, comparator, and falsifier surfaces.


For Falsifiability: Falsifiability K10, reader use is to ask whether the claimed phenomenon would actually reopen under the stated falsifier; if that module-specific question cannot be answered from the promoted theorem, proof sheet, finite semantic case, replay table, comparator row, or claim-boundary section, the source remains background support rather than public theorem evidence.


For Falsifiability: Falsifiability K10, the audit route keeps the complete technical module indexed by the release corpus index and the public evidence package. The science is therefore not discarded, while the PDF avoids becoming a unintegrated technical dossier. A line-level reviewer follows the register; a scientific reader follows this digest and the main chapter.


Falsifiability: Falsifiability K10 reader checkpoint is satisfied only when the main manuscript tells the reader which claim class this module can support, which evidence class limits it, which prior-art or comparator boundary applies, and which failure would reopen the claim. This fourth check is deliberately positive: the module must contribute intelligibility, not merely avoid forbidden wording.


For Falsifiability: Falsifiability K10, the external wording boundary allows public language to say that the route helps organize, test, or constrain a bounded OC claim. It may not say that the route proves full-domain closure, final bounded cross-domain synthesis status, unrestricted comparative claim, or a domain result absent from the current proof and evidence layer. This boundary is part of the source's scientific content, not a marketing caveat.


\paragraph{Falsifiability: Falsifiability K11}

For Falsifiability: Falsifiability K11, the scholarly synthesis contributes falsifier grammar and reopening condition, explains why this module exists, shows how it connects to the current 1.3.3 argument, and prevents stronger inherited language from being promoted without the current proof, replay, comparator, and falsifier surfaces.


For Falsifiability: Falsifiability K11, reader use is to ask whether the claimed phenomenon would actually reopen under the stated falsifier; if that module-specific question cannot be answered from the promoted theorem, proof sheet, finite semantic case, replay table, comparator row, or claim-boundary section, the source remains background support rather than public theorem evidence.


For Falsifiability: Falsifiability K11, the audit route keeps the complete technical module indexed by the release corpus index and the public evidence package. The science is therefore not discarded, while the PDF avoids becoming a unintegrated technical dossier. A line-level reviewer follows the register; a scientific reader follows this digest and the main chapter.


Falsifiability: Falsifiability K11 reader checkpoint is satisfied only when the main manuscript tells the reader which claim class this module can support, which evidence class limits it, which prior-art or comparator boundary applies, and which failure would reopen the claim. This fourth check is deliberately positive: the module must contribute intelligibility, not merely avoid forbidden wording.


For Falsifiability: Falsifiability K11, the external wording boundary allows public language to say that the route helps organize, test, or constrain a bounded OC claim. It may not say that the route proves full-domain closure, final bounded cross-domain synthesis status, unrestricted comparative claim, or a domain result absent from the current proof and evidence layer. This boundary is part of the source's scientific content, not a marketing caveat.


\paragraph{Falsifiability: Falsifiability K12}

For Falsifiability: Falsifiability K12, the scholarly synthesis contributes falsifier grammar and reopening condition, explains why this module exists, shows how it connects to the current 1.3.3 argument, and prevents stronger inherited language from being promoted without the current proof, replay, comparator, and falsifier surfaces.


For Falsifiability: Falsifiability K12, reader use is to ask whether the claimed phenomenon would actually reopen under the stated falsifier; if that module-specific question cannot be answered from the promoted theorem, proof sheet, finite semantic case, replay table, comparator row, or claim-boundary section, the source remains background support rather than public theorem evidence.


For Falsifiability: Falsifiability K12, the audit route keeps the complete technical module indexed by the release corpus index and the public evidence package. The science is therefore not discarded, while the PDF avoids becoming a unintegrated technical dossier. A line-level reviewer follows the register; a scientific reader follows this digest and the main chapter.


Falsifiability: Falsifiability K12 reader checkpoint is satisfied only when the main manuscript tells the reader which claim class this module can support, which evidence class limits it, which prior-art or comparator boundary applies, and which failure would reopen the claim. This fourth check is deliberately positive: the module must contribute intelligibility, not merely avoid forbidden wording.


For Falsifiability: Falsifiability K12, the external wording boundary allows public language to say that the route helps organize, test, or constrain a bounded OC claim. It may not say that the route proves full-domain closure, final bounded cross-domain synthesis status, unrestricted comparative claim, or a domain result absent from the current proof and evidence layer. This boundary is part of the source's scientific content, not a marketing caveat.


\subsubsection{Predictions}

\paragraph{Predictions: Predictions Master}

For Predictions: Predictions Master, the scholarly synthesis contributes prediction-route grammar and empirical promotion boundary, explains why this module exists, shows how it connects to the current 1.3.3 argument, and prevents stronger inherited language from being promoted without the current proof, replay, comparator, and falsifier surfaces.


For Predictions: Predictions Master, reader use is to ask whether the row has formula, snapshot, comparator, uncertainty, residual, negative control, falsifier, and replay hash; if that module-specific question cannot be answered from the promoted theorem, proof sheet, finite semantic case, replay table, comparator row, or claim-boundary section, the source remains background support rather than public theorem evidence.


For Predictions: Predictions Master, the audit route keeps the complete technical module indexed by the release corpus index and the public evidence package. The science is therefore not discarded, while the PDF avoids becoming a unintegrated technical dossier. A line-level reviewer follows the register; a scientific reader follows this digest and the main chapter.


Predictions: Predictions Master reader checkpoint is satisfied only when the main manuscript tells the reader which claim class this module can support, which evidence class limits it, which prior-art or comparator boundary applies, and which failure would reopen the claim. This fourth check is deliberately positive: the module must contribute intelligibility, not merely avoid forbidden wording.


For Predictions: Predictions Master, the external wording boundary allows public language to say that the route helps organize, test, or constrain a bounded OC claim. It may not say that the route proves full-domain closure, final bounded cross-domain synthesis status, unrestricted comparative claim, or a domain result absent from the current proof and evidence layer. This boundary is part of the source's scientific content, not a marketing caveat.


\paragraph{Predictions: Predictions K0}

For Predictions: Predictions K0, the scholarly synthesis contributes prediction-route grammar and empirical promotion boundary, explains why this module exists, shows how it connects to the current 1.3.3 argument, and prevents stronger inherited language from being promoted without the current proof, replay, comparator, and falsifier surfaces.


For Predictions: Predictions K0, reader use is to ask whether the row has formula, snapshot, comparator, uncertainty, residual, negative control, falsifier, and replay hash; if that module-specific question cannot be answered from the promoted theorem, proof sheet, finite semantic case, replay table, comparator row, or claim-boundary section, the source remains background support rather than public theorem evidence.


For Predictions: Predictions K0, the audit route keeps the complete technical module indexed by the release corpus index and the public evidence package. The science is therefore not discarded, while the PDF avoids becoming a unintegrated technical dossier. A line-level reviewer follows the register; a scientific reader follows this digest and the main chapter.


Predictions: Predictions K0 reader checkpoint is satisfied only when the main manuscript tells the reader which claim class this module can support, which evidence class limits it, which prior-art or comparator boundary applies, and which failure would reopen the claim. This fourth check is deliberately positive: the module must contribute intelligibility, not merely avoid forbidden wording.


For Predictions: Predictions K0, the external wording boundary allows public language to say that the route helps organize, test, or constrain a bounded OC claim. It may not say that the route proves full-domain closure, final bounded cross-domain synthesis status, unrestricted comparative claim, or a domain result absent from the current proof and evidence layer. This boundary is part of the source's scientific content, not a marketing caveat.


\paragraph{Predictions: Predictions K1}

For Predictions: Predictions K1, the scholarly synthesis contributes prediction-route grammar and empirical promotion boundary, explains why this module exists, shows how it connects to the current 1.3.3 argument, and prevents stronger inherited language from being promoted without the current proof, replay, comparator, and falsifier surfaces.


For Predictions: Predictions K1, reader use is to ask whether the row has formula, snapshot, comparator, uncertainty, residual, negative control, falsifier, and replay hash; if that module-specific question cannot be answered from the promoted theorem, proof sheet, finite semantic case, replay table, comparator row, or claim-boundary section, the source remains background support rather than public theorem evidence.


For Predictions: Predictions K1, the audit route keeps the complete technical module indexed by the release corpus index and the public evidence package. The science is therefore not discarded, while the PDF avoids becoming a unintegrated technical dossier. A line-level reviewer follows the register; a scientific reader follows this digest and the main chapter.


Predictions: Predictions K1 reader checkpoint is satisfied only when the main manuscript tells the reader which claim class this module can support, which evidence class limits it, which prior-art or comparator boundary applies, and which failure would reopen the claim. This fourth check is deliberately positive: the module must contribute intelligibility, not merely avoid forbidden wording.


For Predictions: Predictions K1, the external wording boundary allows public language to say that the route helps organize, test, or constrain a bounded OC claim. It may not say that the route proves full-domain closure, final bounded cross-domain synthesis status, unrestricted comparative claim, or a domain result absent from the current proof and evidence layer. This boundary is part of the source's scientific content, not a marketing caveat.


\paragraph{Predictions: Predictions K2}

For Predictions: Predictions K2, the scholarly synthesis contributes prediction-route grammar and empirical promotion boundary, explains why this module exists, shows how it connects to the current 1.3.3 argument, and prevents stronger inherited language from being promoted without the current proof, replay, comparator, and falsifier surfaces.


For Predictions: Predictions K2, reader use is to ask whether the row has formula, snapshot, comparator, uncertainty, residual, negative control, falsifier, and replay hash; if that module-specific question cannot be answered from the promoted theorem, proof sheet, finite semantic case, replay table, comparator row, or claim-boundary section, the source remains background support rather than public theorem evidence.


For Predictions: Predictions K2, the audit route keeps the complete technical module indexed by the release corpus index and the public evidence package. The science is therefore not discarded, while the PDF avoids becoming a unintegrated technical dossier. A line-level reviewer follows the register; a scientific reader follows this digest and the main chapter.


Predictions: Predictions K2 reader checkpoint is satisfied only when the main manuscript tells the reader which claim class this module can support, which evidence class limits it, which prior-art or comparator boundary applies, and which failure would reopen the claim. This fourth check is deliberately positive: the module must contribute intelligibility, not merely avoid forbidden wording.


For Predictions: Predictions K2, the external wording boundary allows public language to say that the route helps organize, test, or constrain a bounded OC claim. It may not say that the route proves full-domain closure, final bounded cross-domain synthesis status, unrestricted comparative claim, or a domain result absent from the current proof and evidence layer. This boundary is part of the source's scientific content, not a marketing caveat.


\paragraph{Predictions: Predictions K3}

For Predictions: Predictions K3, the scholarly synthesis contributes prediction-route grammar and empirical promotion boundary, explains why this module exists, shows how it connects to the current 1.3.3 argument, and prevents stronger inherited language from being promoted without the current proof, replay, comparator, and falsifier surfaces.


For Predictions: Predictions K3, reader use is to ask whether the row has formula, snapshot, comparator, uncertainty, residual, negative control, falsifier, and replay hash; if that module-specific question cannot be answered from the promoted theorem, proof sheet, finite semantic case, replay table, comparator row, or claim-boundary section, the source remains background support rather than public theorem evidence.


For Predictions: Predictions K3, the audit route keeps the complete technical module indexed by the release corpus index and the public evidence package. The science is therefore not discarded, while the PDF avoids becoming a unintegrated technical dossier. A line-level reviewer follows the register; a scientific reader follows this digest and the main chapter.


Predictions: Predictions K3 reader checkpoint is satisfied only when the main manuscript tells the reader which claim class this module can support, which evidence class limits it, which prior-art or comparator boundary applies, and which failure would reopen the claim. This fourth check is deliberately positive: the module must contribute intelligibility, not merely avoid forbidden wording.


For Predictions: Predictions K3, the external wording boundary allows public language to say that the route helps organize, test, or constrain a bounded OC claim. It may not say that the route proves full-domain closure, final bounded cross-domain synthesis status, unrestricted comparative claim, or a domain result absent from the current proof and evidence layer. This boundary is part of the source's scientific content, not a marketing caveat.


\paragraph{Predictions: Predictions K4}

For Predictions: Predictions K4, the scholarly synthesis contributes prediction-route grammar and empirical promotion boundary, explains why this module exists, shows how it connects to the current 1.3.3 argument, and prevents stronger inherited language from being promoted without the current proof, replay, comparator, and falsifier surfaces.


For Predictions: Predictions K4, reader use is to ask whether the row has formula, snapshot, comparator, uncertainty, residual, negative control, falsifier, and replay hash; if that module-specific question cannot be answered from the promoted theorem, proof sheet, finite semantic case, replay table, comparator row, or claim-boundary section, the source remains background support rather than public theorem evidence.


For Predictions: Predictions K4, the audit route keeps the complete technical module indexed by the release corpus index and the public evidence package. The science is therefore not discarded, while the PDF avoids becoming a unintegrated technical dossier. A line-level reviewer follows the register; a scientific reader follows this digest and the main chapter.


Predictions: Predictions K4 reader checkpoint is satisfied only when the main manuscript tells the reader which claim class this module can support, which evidence class limits it, which prior-art or comparator boundary applies, and which failure would reopen the claim. This fourth check is deliberately positive: the module must contribute intelligibility, not merely avoid forbidden wording.


For Predictions: Predictions K4, the external wording boundary allows public language to say that the route helps organize, test, or constrain a bounded OC claim. It may not say that the route proves full-domain closure, final bounded cross-domain synthesis status, unrestricted comparative claim, or a domain result absent from the current proof and evidence layer. This boundary is part of the source's scientific content, not a marketing caveat.


\paragraph{Predictions: Predictions K5}

For Predictions: Predictions K5, the scholarly synthesis contributes prediction-route grammar and empirical promotion boundary, explains why this module exists, shows how it connects to the current 1.3.3 argument, and prevents stronger inherited language from being promoted without the current proof, replay, comparator, and falsifier surfaces.


For Predictions: Predictions K5, reader use is to ask whether the row has formula, snapshot, comparator, uncertainty, residual, negative control, falsifier, and replay hash; if that module-specific question cannot be answered from the promoted theorem, proof sheet, finite semantic case, replay table, comparator row, or claim-boundary section, the source remains background support rather than public theorem evidence.


For Predictions: Predictions K5, the audit route keeps the complete technical module indexed by the release corpus index and the public evidence package. The science is therefore not discarded, while the PDF avoids becoming a unintegrated technical dossier. A line-level reviewer follows the register; a scientific reader follows this digest and the main chapter.


Predictions: Predictions K5 reader checkpoint is satisfied only when the main manuscript tells the reader which claim class this module can support, which evidence class limits it, which prior-art or comparator boundary applies, and which failure would reopen the claim. This fourth check is deliberately positive: the module must contribute intelligibility, not merely avoid forbidden wording.


For Predictions: Predictions K5, the external wording boundary allows public language to say that the route helps organize, test, or constrain a bounded OC claim. It may not say that the route proves full-domain closure, final bounded cross-domain synthesis status, unrestricted comparative claim, or a domain result absent from the current proof and evidence layer. This boundary is part of the source's scientific content, not a marketing caveat.


\paragraph{Predictions: Predictions K6}

For Predictions: Predictions K6, the scholarly synthesis contributes prediction-route grammar and empirical promotion boundary, explains why this module exists, shows how it connects to the current 1.3.3 argument, and prevents stronger inherited language from being promoted without the current proof, replay, comparator, and falsifier surfaces.


For Predictions: Predictions K6, reader use is to ask whether the row has formula, snapshot, comparator, uncertainty, residual, negative control, falsifier, and replay hash; if that module-specific question cannot be answered from the promoted theorem, proof sheet, finite semantic case, replay table, comparator row, or claim-boundary section, the source remains background support rather than public theorem evidence.


For Predictions: Predictions K6, the audit route keeps the complete technical module indexed by the release corpus index and the public evidence package. The science is therefore not discarded, while the PDF avoids becoming a unintegrated technical dossier. A line-level reviewer follows the register; a scientific reader follows this digest and the main chapter.


Predictions: Predictions K6 reader checkpoint is satisfied only when the main manuscript tells the reader which claim class this module can support, which evidence class limits it, which prior-art or comparator boundary applies, and which failure would reopen the claim. This fourth check is deliberately positive: the module must contribute intelligibility, not merely avoid forbidden wording.


For Predictions: Predictions K6, the external wording boundary allows public language to say that the route helps organize, test, or constrain a bounded OC claim. It may not say that the route proves full-domain closure, final bounded cross-domain synthesis status, unrestricted comparative claim, or a domain result absent from the current proof and evidence layer. This boundary is part of the source's scientific content, not a marketing caveat.


Integration checkpoint 36. The preceding modules teach inspection logic; they do not inflate claim strength. The release remains bounded by the current proof and evidence surface.


\paragraph{Predictions: Predictions K7}

For Predictions: Predictions K7, the scholarly synthesis contributes prediction-route grammar and empirical promotion boundary, explains why this module exists, shows how it connects to the current 1.3.3 argument, and prevents stronger inherited language from being promoted without the current proof, replay, comparator, and falsifier surfaces.


For Predictions: Predictions K7, reader use is to ask whether the row has formula, snapshot, comparator, uncertainty, residual, negative control, falsifier, and replay hash; if that module-specific question cannot be answered from the promoted theorem, proof sheet, finite semantic case, replay table, comparator row, or claim-boundary section, the source remains background support rather than public theorem evidence.


For Predictions: Predictions K7, the audit route keeps the complete technical module indexed by the release corpus index and the public evidence package. The science is therefore not discarded, while the PDF avoids becoming a unintegrated technical dossier. A line-level reviewer follows the register; a scientific reader follows this digest and the main chapter.


Predictions: Predictions K7 reader checkpoint is satisfied only when the main manuscript tells the reader which claim class this module can support, which evidence class limits it, which prior-art or comparator boundary applies, and which failure would reopen the claim. This fourth check is deliberately positive: the module must contribute intelligibility, not merely avoid forbidden wording.


For Predictions: Predictions K7, the external wording boundary allows public language to say that the route helps organize, test, or constrain a bounded OC claim. It may not say that the route proves full-domain closure, final bounded cross-domain synthesis status, unrestricted comparative claim, or a domain result absent from the current proof and evidence layer. This boundary is part of the source's scientific content, not a marketing caveat.


\paragraph{Predictions: Predictions K8}

For Predictions: Predictions K8, the scholarly synthesis contributes prediction-route grammar and empirical promotion boundary, explains why this module exists, shows how it connects to the current 1.3.3 argument, and prevents stronger inherited language from being promoted without the current proof, replay, comparator, and falsifier surfaces.


For Predictions: Predictions K8, reader use is to ask whether the row has formula, snapshot, comparator, uncertainty, residual, negative control, falsifier, and replay hash; if that module-specific question cannot be answered from the promoted theorem, proof sheet, finite semantic case, replay table, comparator row, or claim-boundary section, the source remains background support rather than public theorem evidence.


For Predictions: Predictions K8, the audit route keeps the complete technical module indexed by the release corpus index and the public evidence package. The science is therefore not discarded, while the PDF avoids becoming a unintegrated technical dossier. A line-level reviewer follows the register; a scientific reader follows this digest and the main chapter.


Predictions: Predictions K8 reader checkpoint is satisfied only when the main manuscript tells the reader which claim class this module can support, which evidence class limits it, which prior-art or comparator boundary applies, and which failure would reopen the claim. This fourth check is deliberately positive: the module must contribute intelligibility, not merely avoid forbidden wording.


For Predictions: Predictions K8, the external wording boundary allows public language to say that the route helps organize, test, or constrain a bounded OC claim. It may not say that the route proves full-domain closure, final bounded cross-domain synthesis status, unrestricted comparative claim, or a domain result absent from the current proof and evidence layer. This boundary is part of the source's scientific content, not a marketing caveat.


\paragraph{Predictions: Predictions K9}

For Predictions: Predictions K9, the scholarly synthesis contributes prediction-route grammar and empirical promotion boundary, explains why this module exists, shows how it connects to the current 1.3.3 argument, and prevents stronger inherited language from being promoted without the current proof, replay, comparator, and falsifier surfaces.


For Predictions: Predictions K9, reader use is to ask whether the row has formula, snapshot, comparator, uncertainty, residual, negative control, falsifier, and replay hash; if that module-specific question cannot be answered from the promoted theorem, proof sheet, finite semantic case, replay table, comparator row, or claim-boundary section, the source remains background support rather than public theorem evidence.


For Predictions: Predictions K9, the audit route keeps the complete technical module indexed by the release corpus index and the public evidence package. The science is therefore not discarded, while the PDF avoids becoming a unintegrated technical dossier. A line-level reviewer follows the register; a scientific reader follows this digest and the main chapter.


Predictions: Predictions K9 reader checkpoint is satisfied only when the main manuscript tells the reader which claim class this module can support, which evidence class limits it, which prior-art or comparator boundary applies, and which failure would reopen the claim. This fourth check is deliberately positive: the module must contribute intelligibility, not merely avoid forbidden wording.


For Predictions: Predictions K9, the external wording boundary allows public language to say that the route helps organize, test, or constrain a bounded OC claim. It may not say that the route proves full-domain closure, final bounded cross-domain synthesis status, unrestricted comparative claim, or a domain result absent from the current proof and evidence layer. This boundary is part of the source's scientific content, not a marketing caveat.


\paragraph{Predictions: Predictions K10}

For Predictions: Predictions K10, the scholarly synthesis contributes prediction-route grammar and empirical promotion boundary, explains why this module exists, shows how it connects to the current 1.3.3 argument, and prevents stronger inherited language from being promoted without the current proof, replay, comparator, and falsifier surfaces.


For Predictions: Predictions K10, reader use is to ask whether the row has formula, snapshot, comparator, uncertainty, residual, negative control, falsifier, and replay hash; if that module-specific question cannot be answered from the promoted theorem, proof sheet, finite semantic case, replay table, comparator row, or claim-boundary section, the source remains background support rather than public theorem evidence.


For Predictions: Predictions K10, the audit route keeps the complete technical module indexed by the release corpus index and the public evidence package. The science is therefore not discarded, while the PDF avoids becoming a unintegrated technical dossier. A line-level reviewer follows the register; a scientific reader follows this digest and the main chapter.


Predictions: Predictions K10 reader checkpoint is satisfied only when the main manuscript tells the reader which claim class this module can support, which evidence class limits it, which prior-art or comparator boundary applies, and which failure would reopen the claim. This fourth check is deliberately positive: the module must contribute intelligibility, not merely avoid forbidden wording.


For Predictions: Predictions K10, the external wording boundary allows public language to say that the route helps organize, test, or constrain a bounded OC claim. It may not say that the route proves full-domain closure, final bounded cross-domain synthesis status, unrestricted comparative claim, or a domain result absent from the current proof and evidence layer. This boundary is part of the source's scientific content, not a marketing caveat.


\paragraph{Predictions: Predictions K11}

For Predictions: Predictions K11, the scholarly synthesis contributes prediction-route grammar and empirical promotion boundary, explains why this module exists, shows how it connects to the current 1.3.3 argument, and prevents stronger inherited language from being promoted without the current proof, replay, comparator, and falsifier surfaces.


For Predictions: Predictions K11, reader use is to ask whether the row has formula, snapshot, comparator, uncertainty, residual, negative control, falsifier, and replay hash; if that module-specific question cannot be answered from the promoted theorem, proof sheet, finite semantic case, replay table, comparator row, or claim-boundary section, the source remains background support rather than public theorem evidence.


For Predictions: Predictions K11, the audit route keeps the complete technical module indexed by the release corpus index and the public evidence package. The science is therefore not discarded, while the PDF avoids becoming a unintegrated technical dossier. A line-level reviewer follows the register; a scientific reader follows this digest and the main chapter.


Predictions: Predictions K11 reader checkpoint is satisfied only when the main manuscript tells the reader which claim class this module can support, which evidence class limits it, which prior-art or comparator boundary applies, and which failure would reopen the claim. This fourth check is deliberately positive: the module must contribute intelligibility, not merely avoid forbidden wording.


For Predictions: Predictions K11, the external wording boundary allows public language to say that the route helps organize, test, or constrain a bounded OC claim. It may not say that the route proves full-domain closure, final bounded cross-domain synthesis status, unrestricted comparative claim, or a domain result absent from the current proof and evidence layer. This boundary is part of the source's scientific content, not a marketing caveat.


\paragraph{Predictions: Predictions K12}

For Predictions: Predictions K12, the scholarly synthesis contributes prediction-route grammar and empirical promotion boundary, explains why this module exists, shows how it connects to the current 1.3.3 argument, and prevents stronger inherited language from being promoted without the current proof, replay, comparator, and falsifier surfaces.


For Predictions: Predictions K12, reader use is to ask whether the row has formula, snapshot, comparator, uncertainty, residual, negative control, falsifier, and replay hash; if that module-specific question cannot be answered from the promoted theorem, proof sheet, finite semantic case, replay table, comparator row, or claim-boundary section, the source remains background support rather than public theorem evidence.


For Predictions: Predictions K12, the audit route keeps the complete technical module indexed by the release corpus index and the public evidence package. The science is therefore not discarded, while the PDF avoids becoming a unintegrated technical dossier. A line-level reviewer follows the register; a scientific reader follows this digest and the main chapter.


Predictions: Predictions K12 reader checkpoint is satisfied only when the main manuscript tells the reader which claim class this module can support, which evidence class limits it, which prior-art or comparator boundary applies, and which failure would reopen the claim. This fourth check is deliberately positive: the module must contribute intelligibility, not merely avoid forbidden wording.


For Predictions: Predictions K12, the external wording boundary allows public language to say that the route helps organize, test, or constrain a bounded OC claim. It may not say that the route proves full-domain closure, final bounded cross-domain synthesis status, unrestricted comparative claim, or a domain result absent from the current proof and evidence layer. This boundary is part of the source's scientific content, not a marketing caveat.
